% !PRINTED 751
\weberpage{751}

%<<CONT
This number $\eta$, which after all likewise belongs to the interval $(0, 1)$, can agree with no number of the sequence $\Omega$.

The formation of $\eta$ can be generalised still further by choosing the initial $\beta_\nu$ arbitrarily through any prescribed finite index, and imposing only thereafter the law that $\beta_\nu$ differ from $\alpha_\nu^{(\nu)}$.

Since it has now been proved that the real algebraic numbers form an enumerable set, and that in every interval there are numbers which do not belong to a given enumerable set, there follows from this quite immediately:

\begin{center}
\gesperrt{In every real interval there are transcendental numbers.}
\end{center}

\begin{center}
\S.\,205.

\gesperrt{Transcendence of the number $e$.}
\end{center}

It is now a far more difficult task to decide, of a definitely given number, whether it is algebraic or transcendental. Here the interest has concentrated chiefly on the two numbers that occur so frequently in analysis: $e$, that is, the base of the natural system of logarithms, and the \gesperrt{Ludolphian} number $\pi$, the ratio of the circumference of the circle to the diameter.

For the number $e$ the question was decided by \gesperrt{Hermite} in a celebrated memoir, which became the foundation for the later investigations on the number $\pi$$^{1)}$. The decision for the number $\pi$, however, which was quite particularly interesting on account of its connection with the long-celebrated problem of the quadrature of the circle, still presented insuperable difficulties for a long time. Finally the demonstration was furnished by \gesperrt{Lindemann} that $\pi$ too belongs to the transcendental numbers. But the proof given by \gesperrt{Lindemann} at first still presented great difficulties to the understanding, which through later investigations by \gesperrt{Weierstrass}, \gesperrt{Hilbert}, \gesperrt{Hurwitz} and \gesperrt{Gordan}$^{2)}$

\noindent\rule{4cm}{0.4pt}

$^{1)}$ \gesperrt{Hermite}, Sur la fonction exponentielle. Comptes rendus, T.\,LXXVII, 1873.

$^{2)}$ \gesperrt{Lindemann}, Ueber die Zahl $\pi$. Mathem. Annalen, Bd. 20, 1882. \gesperrt{Weierstrass}, Zu \gesperrt{Lindemann}'s Abhandlung \glqq Ueber die \gesperrt{Ludolph}'sche
%CONT>>
